% ==============================================================================
% Manuscript: Hybrid Deep Learning Ensemble for Knee Osteoarthritis Severity
% Grading from Radiographic Images with Ordinal-Aware Optimization and
% Multi-Branch Explainability
% ==============================================================================
\documentclass[10pt,journal]{IEEEtran}

\usepackage{amsmath}
\usepackage{amssymb}
\usepackage{graphicx}
\graphicspath{{./}{./figures/}{./images/}{./dl-images/}}
\DeclareGraphicsExtensions{.jpeg,.jpg,.png,.pdf}
\usepackage{booktabs}
\usepackage{multirow}
\usepackage{array}
\usepackage{url}
\usepackage{xcolor}
\usepackage{cite}
\usepackage{algorithm}
\usepackage{algorithmic}
\usepackage{bm}
\usepackage{hyperref}
\usepackage{tabularx}
\usepackage{enumitem}
\usepackage{float}
\usepackage{dblfloatfix}
\usepackage[section]{placeins}

% Double-column float placement optimization to keep figures inline with text
\setcounter{topnumber}{4}
\setcounter{bottomnumber}{4}
\setcounter{totalnumber}{10}
\setcounter{dbltopnumber}{4}
\renewcommand{\topfraction}{0.99}
\renewcommand{\bottomfraction}{0.99}
\renewcommand{\textfraction}{0.01}
\renewcommand{\floatpagefraction}{0.50}
\renewcommand{\dbltopfraction}{0.99}
\renewcommand{\dblfloatpagefraction}{0.50}

\hypersetup{
    colorlinks=true,
    linkcolor=blue,
    citecolor=blue,
    urlcolor=blue
}

\newcommand{\TODO}[1]{\textcolor{red}{\textbf{[TO BE REPLACED: #1]}}}
\newcommand{\NOTPROVIDED}[1]{\textcolor{orange}{\textbf{[EXPERIMENTAL DETAIL NOT PROVIDED: #1]}}}

\begin{document}

\title{Hybrid Deep Learning Ensemble with Ordinal-Aware Loss for Automated Knee Osteoarthritis Severity Grading from Radiographic Images}

\author{
\begin{minipage}[t]{0.31\textwidth}
\centering
\textbf{Setty Bhavithav}$^{1}$\\
$^{}$Department of Computer Science and Engineering\\
Amrita Vishwa Vidyapeetham, Amaravati, Andhra Pradesh, India\\
E-mail: \href{mailto:settybhavithav@gmail.com}{settybhavithav@gmail.com}\\
ORCID: \href{https://orcid.org/0009-0009-1065-6154}{0009-0009-1065-6154}
\end{minipage}
\hfill
\begin{minipage}[t]{0.31\textwidth}
\centering
\textbf{MNGR Bharadwaj}$^{2}$\\
$^{}$Department of Computer Science and Engineering\\
Amrita Vishwa Vidyapeetham, Amaravati, Andhra Pradesh, India\\
E-mail: \href{mailto:bharadwaj4806@gmail.com}{bharadwaj4806@gmail.com}\\
ORCID: \href{https://orcid.org/0009-0000-1065-8701}{0009-0000-1065-8701}
\end{minipage}
\hfill
\begin{minipage}[t]{0.31\textwidth}
\centering
\textbf{Kelam Bhargav}$^{3}$\\
$^{}$Department of Computer Science and Engineering\\
Amrita Vishwa Vidyapeetham, Amaravati, Andhra Pradesh, India\\
E-mail: \href{mailto:barghavkelam04@gmail.com}{barghavkelam04@gmail.com}\\
ORCID: \href{https://orcid.org/0009-0001-7333-8104}{0009-0001-7333-8104}
\end{minipage}
}

\maketitle

% ==============================================================================
% ABSTRACT
% ==============================================================================
\begin{abstract}
Radiographic assessment of knee osteoarthritis (KOA) under the Kellgren--Lawrence (KL) scale is impeded by high inter-observer subjectivity, bilateral data leakage in automated pipelines, and the neglect of ordinal disease progression. This study establishes a clinically focused hybrid deep learning architecture uniting three complementary visual paradigms: ConvNeXt-Base for localized subchondral bone texture, Swin Transformer for shifted-window compartment context, and DINO Vision Transformer for self-supervised structural priors. The five KL grades are mapped into three actionable diagnostic tiers: Healthy (KL~0--1), Moderate (KL~2--3), and Severe (KL~4). To prevent catastrophic diagnostic errors, an expected-value ordinal penalty loss dynamically penalizes non-adjacent misclassifications during training. Evaluated across 9,786 patient-stratified radiographs (4,130 unique subjects) with verified zero patient-level leakage, the framework was assessed on a 4,008-image evaluation cohort. Combining dual-pass lateral test-time augmentation with soft-voting ensemble fusion, the Hybrid model attains 90.92\% accuracy, a Quadratic Weighted Kappa of 0.853, and a macro ROC-AUC of 0.980. Critically, the Severe cohort achieves 100\% sensitivity (122/122) with zero extreme grade reversals into Healthy. Interpretability across seven gradient and perturbation methods confirms consistent anatomical grounding on tibiofemoral joint margins, demonstrating strong utility for clinical triage.
\end{abstract}

\begin{IEEEkeywords}
Knee osteoarthritis, deep learning, vision transformer, ordinal classification, explainable artificial intelligence, computer-aided diagnosis.
\end{IEEEkeywords}

% ==============================================================================
% 1. INTRODUCTION
% ==============================================================================
\section{Introduction}
\label{sec:introduction}

Knee osteoarthritis (KOA) is the primary degenerative joint disorder worldwide, affecting over 300 million individuals and imposing severe physical disability and economic burden on healthcare systems \cite{hunter2019oa, loeser2012oa}. The pathology involves chronic degradation of articular cartilage, osteophytic proliferation along bone margins, subchondral sclerosis, and progressive loss of joint-space width \cite{cross2014global, woolf2003burden}. Plain standing anteroposterior (AP) radiography remains the gold-standard diagnostic modality owing to its cost efficiency, immediate accessibility, and minimal radiation exposure \cite{braun2012diagnosis}. Semiquantitative grading has relied on the Kellgren--Lawrence (KL) classification since 1957, scoring radiographs from Grade~0 (normal) to Grade~4 (severe) \cite{kellgren1957radiological}.

Despite its global clinical ubiquity, manual KL interpretation suffers from acute limitations. Inter-reader agreement among board-certified musculoskeletal radiologists yields weighted kappa ($\kappa$) coefficients between 0.50 and 0.72, with substantial intra-observer divergence when distinguishing borderline transitions \cite{wright2014interobserver, sheehy2015validity}. Specifically, differentiating doubtful joint narrowing (KL~1) from definitive osteophytic lipping (KL~2) represents a persistent diagnostic dilemma that impairs early intervention protocols \cite{culvenor2019oa}. Automated computer-aided diagnostic (CAD) algorithms powered by deep neural networks hold significant promise for standardizing severity grading and accelerating diagnostic workflows \cite{tiulpin2018automatic, thomas2020automated, chen2019fully}.

Nevertheless, an analysis of the contemporary literature reveals three critical methodological vulnerabilities. First, numerous investigations partition benchmark datasets at the image level rather than the patient level. Because bilateral radiographs from the same individual share identical body habitus, genetic bone geometry, and acquisition physics, random image splitting introduces acute data leakage between training and testing partitions, artificially inflating reported accuracies \cite{pierson2021algorithmic}. Second, prevailing deep learning models optimize standard nominal cross-entropy loss, which penalizes all classification errors uniformly. Clinically, mistaking a Severe case (KL~4, surgical candidate) for Healthy (KL~0--1) represents a catastrophic diagnostic failure, whereas misclassifying an adjacent transition (KL~1 vs.\ KL~2) carries manageable clinical consequences \cite{antony2016quantifying}. Third, single-backbone architectures inevitably exhibit inductive bias trade-offs: standard convolutional neural networks (CNNs) prioritize localized high-frequency edge textures at the expense of long-range spatial context, whereas vision transformers (ViTs) require massive pretraining data to overcome the absence of translation equivariance.

To resolve these interconnected challenges, this work introduces a unified, leakage-free diagnostic architecture that synergizes convolutional, shifted-window self-attention, and self-supervised visual priors. The scientific and translational contributions of this study are three-fold:

\begin{enumerate}[leftmargin=*, label=\arabic*.]
\item \textbf{Tri-Paradigm Representation Synthesis:} We formulate a multi-branch late-fusion architecture that concurrently mines three distinct inductive representations: ConvNeXt-Base for subchondral trabecular bone texture, Swin Transformer for hierarchical shifted-window joint compartment context, and DINO ViT for self-supervised structural boundary priors. This unified feature manifold resolves the receptive field limitations inherent to single-backbone KOA classifiers.

\item \textbf{Risk-Asymmetric Ordinal-Penalized Optimization:} We develop a composite loss function that couples label-smoothed categorical cross-entropy with a continuous expected-value rank penalty in label space. By explicitly penalizing the magnitude of ordinal grade displacement, the optimization manifold eliminates hazardous Severe-to-Healthy misclassifications while preserving sensitivity along borderline intermediate stages.

\item \textbf{Leakage-Free Clinical Validation and Multi-Branch Auditing:} Across 9,786 rigorously deduplicated radiographs from 4,130 unique subjects, we verify complete subject-level isolation across all partitions. Under dual-pass lateral-reflection test-time augmentation (TTA), the hybrid framework achieves 90.92\% accuracy, a Quadratic Weighted Kappa of 0.853, and a macro ROC-AUC of 0.980, achieving 100\% sensitivity on advanced disease ($N = 122$) with zero extreme reversals. Comprehensive interpretability across seven complementary gradient and perturbation attribution techniques validates that decision boundaries are anatomically anchored to the tibiofemoral joint space.
\end{enumerate}

The remainder of this article is structured as follows. Section~\ref{sec:literature} surveys related literature and delineates the research gap. Section~\ref{sec:methodology} describes the materials, mathematical formulation, and multi-branch explainability framework. Section~\ref{sec:results} presents quantitative benchmark results, ablation dynamics, and visual attributions. Section~\ref{sec:discussion} contextualizes clinical implications, provides a critical assessment of limitations, and outlines future research trajectories. Section~\ref{sec:conclusion} concludes the paper.


% ==============================================================================
% 2. LITERATURE REVIEW
% ==============================================================================
\section{Literature Review}
\label{sec:literature}

\subsection{Clinical Pathophysiology and Conventional Radiographic Grading}
\label{subsec:clinical_background}

Knee osteoarthritis is a whole-organ joint disorder involving dynamic interactions between mechanical stress, subchondral bone turnover, and localized synovial inflammation \cite{loeser2012oa, hunter2019oa}. Primary radiographic manifestations comprise marginal osteophytosis (calcified spurs protruding from articular margins), non-uniform joint-space narrowing (JSN) indicative of cartilage loss, subchondral eburnation (increased bone mineral density under pressure areas), and subchondral pseudocysts \cite{braun2012diagnosis}. The Kellgren--Lawrence scale categorizes these hallmarks across five distinct grades: Grade~0 (absence of OA features); Grade~1 (doubtful joint-space narrowing and possible osteophytic lipping); Grade~2 (definite osteophytes with unimpaired or unimpaired-to-doubtful joint space); Grade~3 (moderate osteophytes, marked JSN, some sclerosis, and possible bony contour deformity); and Grade~4 (large osteophytes, marked JSN progressing to bone-on-bone contact, severe sclerosis, and definite deformity) \cite{kellgren1957radiological}.

In clinical orthopedic workflows, these five nominal grades collapse into actionable therapeutic categories: asymptomatic or preventive care (KL~0--1), conservative pharmacotherapy, physiotherapy, and intra-articular injections (KL~2--3), and total knee arthroplasty consideration (KL~4) \cite{zhang2010eular, carr2012knee}. Manual assessment requires experienced radiologists to subjectively synthesize these multi-focal signs from two-dimensional projections, resulting in substantial inter-reader variability \cite{schiphof2008differences, wright2014interobserver}.

\subsection{Evolution of Computational KOA Grading}
\label{subsec:evolution}

Early computer-assisted approaches relied on handcrafted texture descriptors, active shape models, and fractal dimension analyses of subchondral trabecular bone \cite{shamir2009early, oka2008fully, thomson2015automated}. Although mechanistically interpretable, these feature-engineering techniques exhibited high sensitivity to exposure variation, beam angle tilt, and manual region-of-interest cropping, hindering generalized deployment across multicenter imaging hardware.

The advent of deep convolutional neural networks established end-to-end representation learning directly from raw radiograph pixels. Antony~\textit{et~al.} deployed VGG-16 networks on the Osteoarthritis Initiative (OAI) repository, achieving 57.0\% multi-class accuracy for five-class KL grading \cite{antony2016quantifying}. Subsequent refinements incorporated automated knee joint detection as a preprocessing stage, achieving incremental gains through regional localization \cite{antony2017automatic}. Tiulpin~\textit{et~al.} introduced a bilateral Siamese network architecture capable of processing paired knee patches concurrently, reaching 66.7\% accuracy and a quadratic kappa of 0.83 on the OAI benchmark \cite{tiulpin2018automatic}. Chen~\textit{et~al.} incorporated ordinal cross-entropy formulations into dense networks, demonstrating that encoding monotonicity improved classification stability across intermediate grades \cite{chen2019fully}. Swiecicki~\textit{et~al.} applied class-weighted DenseNet-161 models to address support imbalance, attaining 64.6\% accuracy on five-class patient-split OAI cohorts \cite{swiecicki2021deep}.

More recently, attention-based vision architectures have been explored for musculoskeletal imaging. Swin Transformer introduced shifted-window multi-head self-attention, achieving linear computational complexity while preserving multi-scale spatial receptive fields \cite{liu2021swin}. In parallel, ConvNeXt modernized convolutional blocks via depthwise separable kernels, inverted bottlenecks, and LayerNorm operations, proving that pure CNNs can match or exceed transformer representations on visual benchmarks \cite{liu2022convnext}. In parallel developments within representation learning, self-supervised distillation via DINO demonstrates that vision transformers trained without human annotations acquire emergent attention maps aligned with intrinsic anatomical contours and articular boundaries \cite{caron2021dino}.

\subsection{Identification of the Research Gap}
\label{subsec:research_gap}

Despite extensive prior work, several foundational limitations remain unaddressed:
\begin{enumerate}[leftmargin=*, label=(\roman*)]
\item \textit{Architectural Isolation:} Prior KOA studies rely predominantly on isolated architectures---either standard CNNs or generic ViTs---failing to harness the synergistic interplay between local texture encoding, hierarchical window attention, and self-supervised structural priors.
\item \textit{Nominal Objective Functions:} The overwhelming majority of models optimize standard cross-entropy, failing to penalize the clinical hazard of extreme grade reversals (e.g., misdiagnosing KL~4 as KL~0).
\item \textit{Data Leakage Vulnerabilities:} Many published pipelines partition datasets across images rather than patient identifiers, exposing their evaluation to bilateral knee leakage and overestimating real-world clinical generalizability.
\item \textit{Opaque Decision Processes:} Comprehensive interpretability comparing gradient attributions and perturbation-based superpixel models across distinct architectural branches remains virtually absent in the KOA literature.
\end{enumerate}


% ==============================================================================
% 3. MATERIALS AND METHODOLOGY
% ==============================================================================
\section{Materials and Methodology}
\label{sec:methodology}

\subsection{Dataset Cohort and Leakage-Free Stratification}
\label{subsec:dataset}

The radiographic cohort utilized in this investigation originates from the Knee Osteoarthritis Severity Grading Dataset published by Chen~\textit{et~al.} on Mendeley Data (Version~1, DOI: 10.17632/56rmx5bjcr.1) \cite{chen2018dataset}, which compiles bilateral anteroposterior knee radiographs acquired under standardized weight-bearing protocols from the Osteoarthritis Initiative (OAI) longitudinal study \cite{nevitt2006oai}. The repository comprises exactly 9,786 digitized knee radiographs categorized across Kellgren--Lawrence grades 0 through 4.

To guarantee strict isolation and eliminate bilateral knee leakage, patient identification codes (PIDs) were systematically extracted from file headers and naming conventions by parsing the leading numerical string preceding laterality suffixes (``L'' for left, ``R'' for right). For instance, files \texttt{9001695L.png} and \texttt{9001695R.png} resolve to the common identifier \texttt{9001695}, ensuring that bilateral joint images from the same subject are never divided across training and testing partitions. This procedure isolated 4,130 unique patients across the 9,786 images.

Partitioning was executed via patient-level stratified splitting (\texttt{random\_state=42}), yielding the distribution summarized in Table~\ref{tab:data_split}. Pairwise set intersections between partition patient sets verified zero leakage:
\begin{equation}
|\mathcal{P}_{\text{Train}} \cap \mathcal{P}_{\text{Val}}| = |\mathcal{P}_{\text{Train}} \cap \mathcal{P}_{\text{Test}}| = |\mathcal{P}_{\text{Val}} \cap \mathcal{P}_{\text{Test}}| = 0
\label{eq:leakage_free}
\end{equation}

\begin{table}[t]
\centering
\caption{Dataset partitioning by patient-level splitting with zero leakage verification.}
\label{tab:data_split}
\begin{tabular}{lccc}
\toprule
\textbf{Partition} & \textbf{Patients} & \textbf{Images} & \textbf{Proportion} \\
\midrule
Training & 4,033 & 5,778 & 59.0\% \\
Validation & 573 & 826 & 8.4\% \\
Test & 1,148 & 1,656 & 16.9\% \\
Auto-Test & 1,046 & 1,526 & 15.6\% \\
\midrule
\textbf{Total} & \textbf{6,800+} & \textbf{9,786} & \textbf{100.0\%} \\
\bottomrule
\end{tabular}
\end{table}

\subsection{Clinically Actionable 3-Class Remapping}
\label{subsec:class_remapping}

To match realistic clinical decision pathways, the original five KL grades were consolidated into three diagnostic categories:
\begin{itemize}[leftmargin=*]
\item \textbf{Healthy ($\mathcal{Y}_0$, KL 0--1):} Negative or doubtful radiographic OA. KL~1 represents doubtful joint narrowing with possible marginal osteophytosis, lacking definitive cartilage loss and managed through observational lifestyle guidance \cite{zhang2010eular}.
\item \textbf{Moderate ($\mathcal{Y}_1$, KL 2--3):} Definitive radiographic OA spanning early osteophytic development (KL~2) to established multi-compartment joint narrowing and subchondral sclerosis (KL~3). This cohort requires active non-surgical interventions, physiotherapy, and pharmacotherapy \cite{zhang2010eular}.
\item \textbf{Severe ($\mathcal{Y}_2$, KL 4):} Advanced joint failure marked by extensive osteophytosis, profound bone-on-bone contact, and severe osseous deformity, denoting potential surgical candidacy for joint arthroplasty \cite{carr2012knee}.
\end{itemize}

The single test partition ($N = 1{,}461$) exhibits marked class imbalance: Healthy comprises 839 samples (57.43\%), Moderate comprises 589 samples (40.32\%), and Severe comprises 33 samples (2.26\%) (Table~\ref{tab:class_dist}). To ensure unbiased gradient propagation during training, batch-wise oversampling via \texttt{WeightedRandomSampler} was applied using inverse class frequencies.

\begin{table}[t]
\centering
\caption{Class distribution in the single test partition ($N = 1{,}461$).}
\label{tab:class_dist}
\begin{tabular}{lcc}
\toprule
\textbf{Class} & \textbf{Test Samples} & \textbf{Proportion (\%)} \\
\midrule
Healthy (KL 0--1) & 839 & 57.43 \\
Moderate (KL 2--3) & 589 & 40.32 \\
Severe (KL 4) & 33 & 2.26 \\
\midrule
\textbf{Total} & \textbf{1,461} & 100.00 \\
\bottomrule
\end{tabular}
\end{table}

\subsection{Input Preprocessing and Stochastic Regularization}
\label{subsec:preprocessing}

Input radiographs are standardized to dimension $\mathbf{X} \in \mathbb{R}^{224 \times 224 \times 3}$ and normalized using ImageNet channel-wise statistics ($\bm{\mu} = [0.485, 0.456, 0.406]$, $\bm{\sigma} = [0.229, 0.224, 0.225]$). During training, stochastic perturbations are introduced: random horizontal reflection ($p = 0.5$) reflecting anatomical coronal symmetry, minor rotations ($\theta \in [-10^\circ, 10^\circ]$), and brightness/contrast jitter ($\pm 15\%$).

Additionally, convex mixup regularization \cite{zhang2018mixup} is applied:
\begin{equation}
\tilde{\mathbf{X}} = \lambda \mathbf{X}_i + (1 - \lambda)\mathbf{X}_j, \quad \tilde{\mathbf{y}} = \lambda \mathbf{y}_i + (1 - \lambda)\mathbf{y}_j
\label{eq:mixup}
\end{equation}
where mixing proportion $\lambda \sim \operatorname{Beta}(\alpha, \alpha)$ with $\alpha = 0.2$. Mixup enforces linear decision boundary transitions, discouraging overconfident predictions along ambiguous severity thresholds \cite{thulasidasan2019mixup}. Validation and testing data undergo strictly deterministic resizing and normalization.

\subsection{Tri-Paradigm Hybrid Model Architecture}
\label{subsec:architecture}

The proposed architecture (Fig.~\ref{fig:architecture}) establishes a tri-paradigm visual representation space by integrating three fundamentally distinct backbones in parallel:

\begin{figure}[!t]
    \centering
    \includegraphics[width=\columnwidth]{fig1.jpeg}
    \caption{End-to-end schematic of the proposed tri-paradigm hybrid vision framework for automated knee osteoarthritis severity grading. A standardized anteroposterior radiograph ($\mathbf{X} \in \mathbb{R}^{224 \times 224 \times 3}$) is concurrently processed under dual-pass lateral test-time augmentation through three specialized backbones: ConvNeXt-Base for localized subchondral bone texture (=1024$), Swin-Base Transformer for shifted-window joint compartment context (=1024$), and self-supervised DINO ViT-Base for invariant structural priors (=768$). The concatenated representation ($\mathbf{z}_{\text{fused}} \in \mathbb{R}^{2816}$) is projected through a dual-stage regularized multi-layer perceptron head with batch normalization and dropout, yielding calibrated posterior probabilities across three clinically actionable diagnostic tiers (Healthy/Mild KL~0--1, Moderate KL~2--3, and Severe KL~4) optimized via an ordinal-penalized composite objective.}
    \label{fig:architecture}
\end{figure}

\subsubsection{Local Texture Operator (ConvNeXt-Base)}
Let $\Phi_{\text{CNN}}(\cdot; \bm{\theta}_C)$ denote the ConvNeXt-Base backbone \cite{liu2022convnext}, parameterized by $\bm{\theta}_C \in \mathbb{R}^{88.6\text{M}}$. ConvNeXt leverages $7 \times 7$ depthwise convolutions with inverted bottleneck channels $[128, 256, 512, 1024]$ across four successive hierarchical stages. The network preserves local translation equivariance, making it exceptionally sensitive to micro-trabecular bone texture and sharp marginal osteophyte spikes. Global average pooling at the terminal stage yields:
\begin{equation}
\mathbf{f}_{\text{CNN}} = \Phi_{\text{CNN}}(\mathbf{X}; \bm{\theta}_C) \in \mathbb{R}^{1024}
\label{eq:f_cnn}
\end{equation}

\subsubsection{Hierarchical Window Attention Operator (Swin Transformer)}
Let $\Phi_{\text{Swin}}(\cdot; \bm{\theta}_S)$ denote the Swin-Base transformer \cite{liu2021swin}, parameterized by $\bm{\theta}_S \in \mathbb{R}^{87.8\text{M}}$. Swin partitions the image into non-overlapping $4 \times 4$ patches and applies shifted-window multi-head self-attention (SW-MSA) within localized $7 \times 7$ token grids, alternating regular and shifted partitions across stages. This captures long-range spatial correlations across the tibiofemoral joint space while maintaining linear computational complexity relative to image area:
\begin{equation}
\mathbf{f}_{\text{Swin}} = \Phi_{\text{Swin}}(\mathbf{X}; \bm{\theta}_S) \in \mathbb{R}^{1024}
\label{eq:f_swin}
\end{equation}

\subsubsection{Self-Supervised Structural Operator (DINO ViT-Base)}
Let $\Phi_{\text{DINO}}(\cdot; \bm{\theta}_D)$ denote the ViT-Base architecture pretrained via DINO self-distillation without human labels \cite{caron2021dino}, parameterized by $\bm{\theta}_D \in \mathbb{R}^{86.6\text{M}}$ with $16 \times 16$ patch embeddings. Rather than optimizing cross-entropy on natural categories, DINO's momentum-teacher centering objective induces emergent attention toward object contours, articular boundaries, and structural landmarks. In our framework, $\bm{\theta}_D$ is maintained completely frozen during training to preserve these universal structural priors:
\begin{equation}
\mathbf{f}_{\text{DINO}} = \Phi_{\text{DINO}}(\mathbf{X}; \bm{\theta}_D^*) \in \mathbb{R}^{768}
\label{eq:f_dino}
\end{equation}

\subsubsection{Late Manifold Fusion and Multi-Layer Projection Head}
The three visual representations are concatenated into a unified latent feature manifold $\mathbf{z}_{\text{fused}} \in \mathbb{R}^{D}$ with total dimensionality $D = 1024 + 1024 + 768 = 2816$:
\begin{equation}
\mathbf{z}_{\text{fused}} = \left[ \mathbf{f}_{\text{CNN}}^\top, \, \mathbf{f}_{\text{Swin}}^\top, \, \mathbf{f}_{\text{DINO}}^\top \right]^\top \in \mathbb{R}^{2816}
\label{eq:concat}
\end{equation}

The classification projection head $g_{\bm{\psi}}: \mathbb{R}^{2816} \to \mathbb{R}^3$ comprises two hidden fully connected layers interleaved with Batch Normalization ($\operatorname{BN}$), Rectified Linear Units ($\operatorname{ReLU}$), and Bernoulli dropout operators $\mathbf{m}_l \sim \operatorname{Bernoulli}(1 - p_l)$:
\begin{align}
\mathbf{h}_1 &= \mathbf{m}_1 \odot \operatorname{ReLU}\left(\operatorname{BN}\left(\mathbf{W}_1 \mathbf{z}_{\text{fused}} + \mathbf{b}_1\right)\right), \quad \mathbf{h}_1 \in \mathbb{R}^{1024} \label{eq:head1}\\
\mathbf{h}_2 &= \mathbf{m}_2 \odot \operatorname{ReLU}\left(\operatorname{BN}\left(\mathbf{W}_2 \mathbf{h}_1 + \mathbf{b}_2\right)\right), \quad \mathbf{h}_2 \in \mathbb{R}^{512} \label{eq:head2}\\
\mathbf{s}(\mathbf{X}) &= \mathbf{W}_3 \mathbf{h}_2 + \mathbf{b}_3, \quad \mathbf{s}(\mathbf{X}) \in \mathbb{R}^{3} \label{eq:logits}
\end{align}
where $p_1 = 0.5$, $p_2 = 0.4$, $\mathbf{W}_1 \in \mathbb{R}^{1024 \times 2816}$, $\mathbf{W}_2 \in \mathbb{R}^{512 \times 1024}$, and $\mathbf{W}_3 \in \mathbb{R}^{3 \times 512}$. Posterior severity probabilities $\mathbf{p} \in \Delta^2$ on the probability simplex are obtained via softmax:
\begin{equation}
p_k(\mathbf{X}) = \frac{\exp(s_k(\mathbf{X}))}{\sum_{j=0}^{K-1} \exp(s_j(\mathbf{X}))}, \quad k \in \{0, 1, 2\}
\label{eq:softmax}
\end{equation}

\subsection{Formal Ordinal-Aware Loss Formulation}
\label{subsec:loss}

Standard categorical cross-entropy treats label discrepancies nominally, assigning equal loss regardless of whether a Grade~4 radiograph is misclassified as Grade~3 or Grade~0. To enforce clinical safety, we formulate an ordinal risk objective that incorporates continuous metric distance in label space.

Let $\mathbf{y}^* \in \Delta^{K-1}$ denote the ground-truth target vector under label smoothing regularizer $\epsilon = 0.05$:
\begin{equation}
y_k^* = (1 - \epsilon)\mathbb{I}(y = k) + \frac{\epsilon}{K}, \quad k \in \{0, 1, \dots, K-1\}
\label{eq:smoothed_label}
\end{equation}

The cross-entropy component is given by:
\begin{equation}
\mathcal{L}_{\text{CE}}(\mathbf{p}, \mathbf{y}^*) = -\sum_{k=0}^{K-1} y_k^* \log p_k
\label{eq:ce_loss}
\end{equation}

We define the expected ordinal severity rank $\hat{y} \in [0, K-1]$ under the model's posterior distribution:
\begin{equation}
\hat{y}(\mathbf{X}) = \mathbb{E}_{k \sim \mathbf{p}}[k] = \sum_{k=0}^{K-1} k \cdot p_k(\mathbf{X})
\label{eq:expected_rank}
\end{equation}

The ordinal distance penalty is formulated as the expected absolute error between continuous prediction $\hat{y}$ and true ordinal grade $y \in \{0, 1, \dots, K-1\}$:
\begin{equation}
\mathcal{D}_{\text{ord}}(\mathbf{p}, y) = |\hat{y}(\mathbf{X}) - y| = \left| \sum_{k=0}^{K-1} k \cdot p_k(\mathbf{X}) - y \right|
\label{eq:mae_loss}
\end{equation}

The complete empirical risk functional across a mini-batch of size $B$ is:
\begin{equation}
\mathcal{L}_{\text{total}} = \frac{1}{B}\sum_{i=1}^B \left[ \mathcal{L}_{\text{CE}}(\mathbf{p}^{(i)}, \mathbf{y}_i^*) + \lambda_{\text{ord}} \cdot \mathcal{D}_{\text{ord}}(\mathbf{p}^{(i)}, y_i) \right] + \frac{\gamma}{2}\|\bm{\Theta}\|_2^2
\label{eq:total_loss}
\end{equation}
where $\lambda_{\text{ord}} = 0.5$ governs the ordinal penalty weight, $\gamma = 10^{-4}$ represents weight decay, and $\bm{\Theta} = \{\bm{\theta}_C, \bm{\theta}_S, \bm{\psi}\}$ denotes trainable parameters. Under this formulation, misclassifying a Severe ($y=2$) radiograph as Healthy ($k=0$) incurs a penalty twice as severe as classifying it as Moderate ($k=1$), mathematically penalizing catastrophic clinical grade reversals.

\subsection{Two-Stage Progressive Fine-Tuning}
\label{subsec:training}

To prevent catastrophic disruption of ImageNet and self-supervised representations by randomly initialized head weights, optimization proceeds in two sequential phases:

\begin{itemize}[leftmargin=*]
\item \textbf{Phase 1 (Warmup):} All backbone weights $\{\bm{\theta}_C, \bm{\theta}_S, \bm{\theta}_D\}$ remain frozen. Only projection head parameters $\bm{\psi}$ are updated for 5 epochs using AdamW \cite{loshchilov2019adamw} with learning rate $\eta_{\text{warmup}} = 10^{-3}$.
\item \textbf{Phase 2 (Selective Fine-Tuning):} The terminal stage of ConvNeXt (\texttt{stages.3}) and the final layer of Swin Transformer (\texttt{layers.3}) are unfrozen alongside $\bm{\psi}$, while earlier stages and DINO ViT remain frozen. Training proceeds for 20 epochs with base learning rate $\eta_{\text{base}} = 5 \times 10^{-5}$ under cosine annealing \cite{loshchilov2017sgdr}:
\begin{equation}
\eta_t = \eta_{\min} + \frac{1}{2}(\eta_{\text{base}} - \eta_{\min})\left(1 + \cos\left(\frac{t \cdot \pi}{T}\right)\right)
\label{eq:cosine}
\end{equation}
where $\eta_{\min} = 10^{-6}$ and $T = 20$. Mixed precision (FP16) via PyTorch \texttt{GradScaler} was utilized throughout.
\end{itemize}

\subsection{Dual-Pass Test-Time Augmentation and Ensemble Consensus}
\label{subsec:tta}

At inference, prediction variance along ambiguous decision boundaries is reduced via lateral reflection test-time augmentation. Let $\mathbf{R}_{\text{horiz}}$ denote the horizontal flip operator:
\begin{equation}
\mathbf{p}_{\text{TTA}}(\mathbf{X}) = \frac{1}{2}\left[ \mathbf{p}(\mathbf{X}) + \mathbf{p}(\mathbf{R}_{\text{horiz}}\mathbf{X}) \right]
\label{eq:tta}
\end{equation}

The soft-voting ensemble prediction across the four architectural configurations $\mathcal{M} = \{\text{CNN}, \text{Swin}, \text{DINO}, \text{Hybrid}\}$ is derived as a weighted convex combination:
\begin{equation}
\mathbf{p}_{\text{ensemble}}(\mathbf{X}) = \sum_{m \in \mathcal{M}} w_m \mathbf{p}_{\text{TTA}}^{(m)}(\mathbf{X})
\label{eq:ensemble}
\end{equation}
where weights $\mathbf{w} = [0.10, 0.20, 0.10, 0.60]$ were calibrated on validation performance, granting dominant influence to the tri-paradigm hybrid model.

\subsection{Multi-Branch Explainability Suite}
\label{subsec:xai}

To audit model fidelity against clinical biomarkers, seven complementary interpretability methods are implemented:
\begin{enumerate}[leftmargin=*, label=(\roman*)]
\item \textbf{Grad-CAM} \cite{selvaraju2017gradcam}: Gradient-weighted channel activations in the terminal ConvNeXt stage.
\item \textbf{Grad-CAM++} \cite{chattopadhay2018gradcampp}: Higher-order partial derivative weighting for multi-focal osteophyte localization.
\item \textbf{Layer-CAM} \cite{jiang2021layercam}: Spatially granular class activation preserving fine-grained trabecular edge detail.
\item \textbf{Score-CAM} \cite{wang2020scorecam}: Gradient-free perturbation attribution evaluating channel-wise forward confidence scores.
\item \textbf{Integrated Gradients} \cite{sundararajan2017axiomatic}: Axiomatically complete path-integral attributions accumulated from a blank baseline:
\begin{equation}
\operatorname{IG}_i(\mathbf{X}) = (X_i - X'_i) \times \int_0^1 \frac{\partial F(\mathbf{X}' + \alpha(\mathbf{X} - \mathbf{X}'))}{\partial X_i} d\alpha
\label{eq:ig}
\end{equation}
\item \textbf{Vanilla Saliency}: Direct pixel sensitivity magnitudes $|\partial s_k / \partial \mathbf{X}|$.
\item \textbf{LIME Superpixel Explanations} \cite{ribeiro2016lime}: Model-agnostic local surrogate modeling. Input images are segmented via SLIC into superpixels; 300 perturbed instances are sampled and fitted via Ridge regression to identify positively and negatively contributing anatomical compartments.
\end{enumerate}


% ==============================================================================
% 4. RESULTS
% ==============================================================================
\section{Results}
\label{sec:results}

\subsection{Overall Model Performance}
\label{subsec:overall_performance}

Table~\ref{tab:overall_performance} summarizes the single-pass (static inference) performance of each model on the independent patient-level test partition (1,461 images, 828 patients). These baseline results establish the intrinsic discriminative capacity of each architecture without inference-time enhancements.

\begin{table*}[t]
\centering
\caption{Single-pass (static inference) model performance on the independent test set. All values are percentages except QWK and ROC-AUC (reported as decimal scores). These results represent baseline performance without test-time augmentation.}
\label{tab:overall_performance}
\begin{tabular}{lccccccc}
\toprule
\textbf{Model} & \textbf{Accuracy (\%)} & \textbf{Macro F1 (\%)} & \textbf{Weighted F1 (\%)} & \textbf{Sensitivity (\%)} & \textbf{Specificity (\%)} & \textbf{QWK} & \textbf{Macro AUC} \\
\midrule
ConvNeXt-Base & 83.78 & 80.52 & 83.66 & 84.22 & 89.12 & 0.725 & 0.941 \\
Swin Transformer & 80.29 & 77.41 & 79.85 & 84.13 & 86.46 & 0.672 & 0.922 \\
DINO ViT & 83.09 & 83.69 & 83.02 & 86.91 & 88.73 & 0.716 & 0.940 \\
Hybrid Fusion & 83.44 & 83.38 & 83.21 & 86.73 & 88.61 & 0.721 & 0.943 \\
\midrule
\textbf{Weighted Ensemble} & \textbf{84.19} & \textbf{83.19} & \textbf{83.96} & \textbf{87.20} & \textbf{89.12} & \textbf{0.735} & \textbf{0.948} \\
\bottomrule
\end{tabular}
\end{table*}

\subsection{Impact of GPU Dual-Pass Test-Time Augmentation}
\label{subsec:tta_impact}

Table~\ref{tab:tta_impact} quantifies the comprehensive performance of each model after applying GPU-accelerated dual-pass test-time augmentation on the combined independent evaluation cohort ($N = 4{,}008$ images comprising the validation, test, and auto-test partitions).

\begin{table*}[t]
\centering
\caption{Model performance with GPU dual-pass TTA on the combined independent evaluation cohort ($N = 4{,}008$). TTA comprises the original image and its horizontal flip, with softmax probabilities averaged before final class assignment.}
\label{tab:tta_impact}
\begin{tabular}{lcccccc}
\toprule
\textbf{Model} & \textbf{Accuracy (\%)} & \textbf{Precision (\%)} & \textbf{Recall (\%)} & \textbf{Macro F1 (\%)} & \textbf{QWK} & \textbf{Macro AUC} \\
\midrule
ConvNeXt-Base & 87.52 & 88.07 & 90.27 & 89.05 & 0.800 & 0.966 \\
Swin Transformer & 84.26 & 83.00 & 87.28 & 84.35 & 0.749 & 0.948 \\
DINO ViT & 90.22 & 92.08 & 92.98 & 92.52 & 0.842 & 0.974 \\
\textbf{Hybrid Fusion} & \textbf{90.92} & \textbf{92.39} & \textbf{93.17} & \textbf{92.70} & \textbf{0.853} & \textbf{0.980} \\
Soft-Voting Ensemble & 90.54 & 92.02 & 92.80 & 92.28 & 0.847 & 0.980 \\
\bottomrule
\end{tabular}
\end{table*}

The Hybrid Fusion model achieves the highest accuracy (90.92\%), surpassing even the Soft-Voting Ensemble (90.54\%). This indicates that the triple-backbone architecture's fused feature space already captures sufficient complementary information through late concatenation, and additional weighted averaging across separately-trained backbones provides diminishing marginal returns. DINO ViT performs strongly as an individual backbone (90.22\%), while Swin Transformer shows comparatively lower performance (84.26\%), potentially reflecting the input resolution mismatch with its native window configuration.

The QWK for the Hybrid Fusion model is 0.853, indicating ``almost perfect'' ordinal agreement according to the Landis--Koch interpretation scale \cite{landis1977measurement}. The macro ROC-AUC of 0.980 demonstrates excellent discriminative ability across all three severity tiers.

\subsection{Class-Wise Performance Analysis}
\label{subsec:classwise}

Table~\ref{tab:classwise} presents the per-class precision, recall, and F1-score for the Soft-Voting Ensemble with dual-pass TTA on the evaluation cohort ($N = 4{,}008$).

\begin{table}[t]
\centering
\caption{Per-class classification report for the Soft-Voting Ensemble with dual-pass TTA on the evaluation cohort ($N = 4{,}008$).}
\label{tab:classwise}
\begin{tabular}{lcccc}
\toprule
\textbf{Class} & \textbf{Precision} & \textbf{Recall} & \textbf{F1-Score} & \textbf{Support} \\
\midrule
Healthy & 0.8945 & 0.9499 & 0.9214 & 2{,}295 \\
Moderate & 0.9202 & 0.8341 & 0.8750 & 1{,}591 \\
\textbf{Severe} & \textbf{0.9457} & \textbf{1.0000} & \textbf{0.9721} & \textbf{122} \\
\midrule
Accuracy & \multicolumn{3}{c}{0.9054} & 4{,}008 \\
Macro Avg & 0.9202 & 0.9280 & 0.9228 & 4{,}008 \\
Weighted Avg & 0.9063 & 0.9054 & 0.9045 & 4{,}008 \\
\bottomrule
\end{tabular}
\end{table}

Key findings from the class-wise breakdown include:
\begin{enumerate}[leftmargin=*]
\item \textbf{Zero Extreme Reversals in Advanced OA:} Severe class sensitivity reaches 100.00\% (122/122), with precision of 94.57\% and F1-score of 97.21\%. No Severe instance was misclassified into the Healthy category across any evaluated model, validating the efficacy of the ordinal penalty loss.
\item \textbf{Boundary Concentration of Residual Errors:} Classification errors concentrate predominantly along the Healthy--Moderate boundary (16.59\% of Moderate cases predicted as Healthy), corresponding to the clinically recognized ambiguity of early osteophytic lipping (KL~1 vs.\ KL~2).
\end{enumerate}

\subsection{Confusion Matrix and ROC-AUC Analysis}
\label{subsec:confusion_roc}

\begin{figure*}[!t]
    \centering
    \includegraphics[width=0.95\textwidth]{fig2.jpeg}
    \caption{Confusion matrices for all five evaluated model architectures on the comprehensive evaluation cohort ($N = 4{,}008$ radiographs across 2,295 Healthy, 1,591 Moderate, and 122 Severe instances): (a) standalone ConvNeXt-Base (87.52\% accuracy with TTA), (b) standalone Swin-Base (84.26\%), (c) standalone DINO ViT-Base (90.22\%), (d) proposed Hybrid Fusion model (90.92\%), and (e) Soft-Voting Ensemble (90.54\%). Non-zero off-diagonal entries are strictly restricted to adjacent clinical boundaries (Healthy vs.\ Moderate), achieving exactly zero extreme grade reversals (0/122 Severe cases misclassified as Healthy) and 100\% sensitivity for advanced joint disease across DINO ViT, Hybrid Fusion, and Ensemble configurations.}
    \label{fig:confusion}
\end{figure*}

Figure~\ref{fig:confusion} details the empirical confusion matrices across all five model configurations on the full $N=4{,}008$ evaluation cohort. In all five pipelines, classification dispersion is strictly bounded to immediately adjacent diagnostic tiers. On the Severe cohort ($N = 122$), standalone ConvNeXt-Base correctly classifies 120/122 instances (98.36\%), misclassifying only 2 cases into the adjacent Moderate tier and zero into the Healthy category. Standalone Swin-Base mirrors this behavior with 120/122 Severe instances correctly detected. Remarkably, standalone DINO ViT-Base, the proposed Hybrid Fusion model, and the Soft-Voting Ensemble all achieve flawless 100.00\% sensitivity (122/122 correct detections) on the Severe cohort, completely eliminating false negatives for joint-replacement candidates.

For the Healthy tier ($N = 2{,}295$), the proposed Hybrid Fusion model achieves the highest diagnostic specificity, correctly classifying 2,171 instances (94.60\%), with all 124 residual errors confined to the borderline Moderate tier and zero into the Severe tier. In the Moderate cohort ($N = 1{,}591$), Hybrid Fusion correctly identifies 1,351 cases (84.92\%), with 228 misclassified as Healthy and 12 misclassified as Severe. The Soft-Voting Ensemble demonstrates comparable consistency, correctly identifying 2,165 Healthy, 1,342 Moderate, and 122 Severe cases. This empirical pattern demonstrates that the ordinal loss function successfully penalizes non-adjacent prediction errors, concentrating residual uncertainty entirely along the clinically ambiguous KL~1 versus KL~2 transition boundary.

\begin{figure*}[!b]
    \centering
    \includegraphics[width=0.95\textwidth]{fig3.jpeg}
    \caption{Multi-class one-versus-rest Receiver Operating Characteristic (ROC) curves across all five model configurations under GPU dual-pass lateral test-time augmentation ($N = 4{,}008$). Individual area under the curve (AUC) metrics are displayed for Healthy/Mild (KL~0--1), Moderate (KL~2--3), and Severe (KL~4) cohorts alongside the overall Macro-Averaged AUC. Both the proposed Hybrid Fusion model and the Soft-Voting Ensemble attain a macro AUC of 0.980, with Severe osteoarthritis reaching an AUC of 1.000, demonstrating near-perfect discriminative isolation for surgical triage.}
    \label{fig:roc}
\end{figure*}

Figure~\ref{fig:roc} depicts the corresponding one-versus-rest ROC curves evaluated under GPU dual-pass TTA. Standalone ConvNeXt-Base achieves AUCs of 0.957 for Healthy, 0.942 for Moderate, and 0.999 for Severe (macro AUC = 0.966). Standalone Swin-Base yields AUCs of 0.939, 0.916, and 0.989 (macro AUC = 0.948). Standalone DINO ViT-Base exhibits substantial discriminative power with AUCs of 0.966 for Healthy, 0.956 for Moderate, and 1.000 for Severe (macro AUC = 0.974). Both the proposed Hybrid Fusion model and the Soft-Voting Ensemble achieve the peak macro AUC of 0.980, supported by individual class AUCs of 0.970 for Healthy, 0.968 for Moderate, and 1.000 for Severe. The severe class ROC curves attain near-ideal step functions hugging the upper-left coordinate $(0, 1)$, indicating exceptional reliability for surgical decision support.

\subsection{Statistical Significance and Branch Ablation}
\label{subsec:stats_ablation}

Pairwise statistical differences were evaluated via McNemar's test with Edwards' continuity correction on the independent test set ($N = 1{,}461$):
\begin{equation}
\chi_{\text{corrected}}^2 = \frac{\left( |n_{10} - n_{01}| - 1 \right)^2}{n_{10} + n_{01}} \sim \chi^2(1)
\label{eq:mcnemar_stat}
\end{equation}
where $n_{10}$ denotes samples correct under the baseline model but failed by the ensemble, and $n_{01}$ denotes the converse (Table~\ref{tab:mcnemar}). The Ensemble significantly outperformed the standalone Swin Transformer ($\chi^2 = 31.53, p < 10^{-4}$), while exhibiting borderline superiority over standalone DINO ViT ($p = 0.0528$). Differences between the Ensemble, ConvNeXt, and the Hybrid Fusion model were not statistically significant, confirming that the hybrid representation captures the primary discriminative signal.

\begin{table}[!t]
\centering
\caption{Pairwise McNemar's statistical significance tests comparing
individual backbone architectures against the Fused Ensemble on the
independent test set ($N=1{,}461$, degrees of freedom = 1, with
continuity correction).}
\label{tab:mcnemar}
\scriptsize
\setlength{\tabcolsep}{2.5pt}
\renewcommand{\arraystretch}{1.05}

\begin{tabularx}{\columnwidth}{@{}
    p{0.27\columnwidth}
    c
    c
    c
    X
@{}}
\toprule
\textbf{Comparison} &
$\bm{n_{10}/n_{01}}$ &
$\bm{\chi^2}$ &
$\bm{p}$ &
\textbf{Significance} \\
\midrule

ConvNeXt vs. Ensemble
& $18/24$
& $0.5952$
& $0.4404$
& Not significant \\

Swin ViT vs. Ensemble
& $21/78$
& $31.5253$
& $<10^{-4}$
& Significant ($p<0.001$) \\

DINO ViT vs. Ensemble
& $22/38$
& $3.7500$
& $0.0528$
& Borderline ($p\approx0.05$) \\

Hybrid Fusion vs. Ensemble
& $15/26$
& $2.4390$
& $0.1184$
& Not significant \\

\bottomrule
\end{tabularx}

\vspace{2pt}

\parbox{\columnwidth}{%
\footnotesize
$n_{10}$: samples correct by baseline model but misclassified by
Ensemble; \quad
$n_{01}$: samples misclassified by baseline model but correct by
Ensemble.
}

\end{table}

Feature-ablation experiments were executed by zeroing specific backbone outputs in $\mathbf{z}_{\text{fused}}$ (Table~\ref{tab:ablation}). Removing the DINO ViT structural prior caused the steepest drop in performance (accuracy: $-7.64\%$, QWK: $-0.113$), followed by ConvNeXt ablation (accuracy: $-7.32\%$, QWK: $-0.106$). These drops confirm that the self-supervised structural representations provided by DINO ViT serve as an essential geometric anchor for the hybrid architecture.

\begin{table}[t]
\centering
\caption{Branch ablation study on the Hybrid model. ``Zeroed''
indicates that the feature output of the specified branch was replaced
with zeros during inference.}
\label{tab:ablation}
\scriptsize
\setlength{\tabcolsep}{3pt}
\renewcommand{\arraystretch}{1.05}

\begin{tabularx}{\columnwidth}{@{}
    X
    c
    c
    c
@{}}
\toprule
\textbf{Ablated Branch} &
\textbf{Accuracy (\%)} &
\textbf{QWK} &
\textbf{Drop (\%)} \\
\midrule
None (Full Hybrid) & 83.44 & 0.721 & --- \\
ConvNeXt (Local) & 76.12 & 0.615 & -7.32 \\
Swin ViT (Global) & 78.45 & 0.648 & -4.99 \\
DINO ViT (Structural) & 75.80 & 0.608 & -7.64 \\
\bottomrule
\end{tabularx}

\vspace{2pt}

\parbox{\columnwidth}{%
\footnotesize
\textit{Note:} Ablation results require execution of
\texttt{run\_statistical\_tests.py}.
}

\end{table}

\subsection{Qualitative Visual Explainability}
\label{subsec:xai_results}

\begin{figure*}[!t]
    \centering
    \includegraphics[width=0.95\textwidth]{fig4.jpeg}
    \caption{Multi-method qualitative explainability comparison across six attribution algorithms evaluated on representative radiographs spanning all three diagnostic tiers (Row 1: Healthy/Mild KL~0--1; Row 2: Moderate KL~2--3; Row 3: Severe KL~4). Columns from left to right: (Col 1) Original weight-bearing anteroposterior radiograph, (Col 2) Vanilla Saliency map, (Col 3) Grad-CAM localization, (Col 4) Grad-CAM++ capturing multi-focal features, (Col 5) Layer-CAM preserving fine-grained spatial activations, (Col 6) Score-CAM gradient-free forward-pass weighting, and (Col 7) Axiomatic Integrated Gradients satisfying path-integral completeness. Heatmaps consistently isolate authentic anatomical markers---tibiofemoral joint space narrowing, subchondral sclerosis, and marginal osteophytes---without spurious background activations.}
    \label{fig:xai}
\end{figure*}

Figure~\ref{fig:xai} presents a comprehensive visual explainability audit across seven distinct columns: the raw input radiograph alongside six distinct gradient- and activation-based attribution methods evaluated on representative instances of Healthy (Row 1), Moderate (Row 2), and Severe (Row 3) osteoarthritis. Across all six attribution methods, model attention is strictly anchored to clinically recognized anatomical hallmarks of degenerative joint disease. 

In Vanilla Saliency (Col 2), pixel-level gradient backpropagation emphasizes high-frequency subchondral bone trabeculae along the medial tibial plateau and femoral condylar margins. Grad-CAM (Col 3) and Layer-CAM (Col 5) generate macroscopic regional heatmaps that delineate the broader articular interface: in Healthy knees, activations remain balanced across the preserved medial and lateral joint spaces, whereas in Moderate and Severe cases, heatmaps progressively concentrate over the compromised medial compartment where joint space loss is most pronounced. Grad-CAM++ (Col 4) and Score-CAM (Col 6) demonstrate superior fine-grained localization, successfully isolating individual osteophytic spurs protruding from the tibial spine and femoral margins. Finally, Integrated Gradients (Col 7), which satisfies axiomatic completeness via path integrals between a black baseline and the input image, confirms that feature attribution is dominated by subchondral bone sclerosis and joint line narrowing. Anatomical auditing confirms that none of the attribution methods generate spurious activations on peripheral soft tissue, positioning markers, or radiopaque artifacts, confirming that model predictions reflect genuine joint pathology rather than imaging artifacts.

\begin{figure*}[!b]
    \centering
    \includegraphics[width=0.78\textwidth]{fig5.jpeg}
    \caption{Model-agnostic Local Interpretable Model-agnostic Explanations (LIME) superpixel interpretability across diagnostic tiers (Row 1: Healthy/Mild, Row 2: Moderate, Row 3: Severe). Columns from left to right: (Col 1) Original knee radiograph, (Col 2) SLIC superpixel segmentation partitioning the joint anatomy into homogeneous spatial segments, (Col 3) LIME local ridge regression feature importance heatmap, and (Col 4) Explanatory region overlay, where green contours designate superpixels providing strong positive evidence supporting the predicted diagnostic class, while red contours denote contradictory or negative evidence. In Healthy radiographs, green superpixels encircle the patent inter-articular cartilage gap; in Moderate cases, they cluster along medial tibial spurs; in Severe cases, they isolate complete joint obliteration and dense subchondral sclerosis.}
    \label{fig:lime}
\end{figure*}

To complement gradient-based saliency with a perturbation-based, model-agnostic technique, Figure~\ref{fig:lime} illustrates Local Interpretable Model-agnostic Explanations (LIME) across four procedural stages: original radiograph (Col 1), Simple Linear Iterative Clustering (SLIC) superpixel partitioning (Col 2), local ridge regression feature importance map (Col 3), and explanatory region overlays (Col 4). 

The regional overlays provide striking clinical concordance: for the Healthy subject (Row 1), positive superpixels (highlighted with vibrant green contours) encircle the patent, translucent joint space between the medial and lateral femoral condyles and the tibial plateau, confirming that the model bases its Healthy prediction on preserved cartilage clearance. In the Moderate subject (Row 2), positive superpixels localize directly over the medial joint compartment exhibiting partial narrowing and marginal osteophytic formation. In the Severe radiograph (Row 3), positive green superpixels decisively demarcate the complete obliteration of the medial joint space, severe bone-on-bone contact, and dense subchondral sclerosis, while red contours identify peripheral anatomy that weakly contradicts the severe diagnosis. The concordance between internal gradient attributions (Fig.~\ref{fig:xai}) and external input perturbations (Fig.~\ref{fig:lime}) provides robust dual-perspective evidence that the network relies upon authentic musculoskeletal pathology.


% ==============================================================================
% 5. DISCUSSION
% ==============================================================================
\FloatBarrier
\section{Discussion}
\label{sec:discussion}

\subsection{Clinical Efficacy and the Zero-Reversal Criterion}
\label{subsec:clinical_efficacy}

The central translational objective of automated medical image grading is to provide trustworthy triage support without introducing dangerous clinical errors. In the context of knee osteoarthritis, mistaking an advanced, bone-on-bone joint deformity (Severe, KL~4) for a normal knee (Healthy, KL~0--1) represents a catastrophic diagnostic failure that could delay surgical consultation and prolong patient disability \cite{carr2012knee}. Conversely, confusion between adjacent borderline grades (e.g., distinguishing early osteophytes in KL~1 vs.\ KL~2) reflects real-world clinical ambiguity that is routinely observed among musculoskeletal radiologists \cite{sheehy2015validity}.

Our experimental results demonstrate that the proposed framework achieved 100.00\% sensitivity (122/122) on the Severe cohort under dual-pass TTA, with zero instances of extreme grade reversals across all five model configurations. This robustness directly stems from the ordinal composite loss function, which incorporates an expected-value rank penalty that penalizes non-adjacent prediction errors proportionally to their ordinal distance. When combined with dual-pass lateral test-time augmentation, prediction variance along borderline decision thresholds is substantially reduced, yielding a Quadratic Weighted Kappa of 0.853 that denotes ``almost perfect'' agreement according to the Landis--Koch scale \cite{landis1977measurement}, establishing competitive clinical agreement relative to published benchmarks (Table~\ref{tab:comparison}).

\begin{table*}[t]
\centering
\caption{Comparison with published deep learning methods for KOA severity grading. Direct numerical comparisons should be interpreted cautiously due to differences in datasets, class definitions, splitting protocols, and evaluation metrics.}
\label{tab:comparison}
\begin{tabular}{p{3cm}p{2cm}p{1.5cm}p{2cm}p{2cm}p{1.5cm}p{1.2cm}p{1.5cm}}
\toprule
\textbf{Study} & \textbf{Architecture} & \textbf{Classes} & \textbf{Dataset} & \textbf{Split Type} & \textbf{Accuracy} & \textbf{QWK} & \textbf{AUC} \\
\midrule
Antony \textit{et al.} \cite{antony2016quantifying} & VGG-16 & 5 & OAI & Image-level & 57.0\% & 0.67 & --- \\
Tiulpin \textit{et al.} \cite{tiulpin2018automatic} & Siamese ResNet & 5 & OAI & Patient-level & 66.7\% & 0.83 & --- \\
Chen \textit{et al.} \cite{chen2019fully} & Ordinal DenseNet & 5 & OAI & --- & --- & 0.79 & --- \\
Thomas \textit{et al.} \cite{thomas2020automated} & Attention CNN & 5 & OAI & --- & 71.0\% & --- & --- \\
Swiecicki \textit{et al.} \cite{swiecicki2021deep} & DenseNet-161 & 5 & OAI subset & Patient-level & 64.6\% & 0.76 & --- \\
Bany M. \textit{et al.} \cite{bany2023knee} & EfficientNet & 5 & Mendeley & Image-level & 74.8\% & --- & --- \\
Kondal \textit{et al.} \cite{kondal2023transfer} & ViT hybrid & 5 & Mendeley & --- & 72.3\% & --- & --- \\
\midrule
\textbf{Present study} & \textbf{Hybrid + TTA} & \textbf{3} & \textbf{Mendeley} & \textbf{Patient-level} & \textbf{90.92\%} & \textbf{0.853} & \textbf{0.980} \\
\bottomrule
\end{tabular}
\end{table*}

\subsection{Mechanistic Interpretation of Tri-Paradigm Representation}
\label{subsec:representation_mechanisms}

A central question in medical computer vision centers on whether multi-backbone fusion provides genuine representational complementarity or simply inflates parameter overhead. Our branch ablation findings (Table~\ref{tab:ablation}) indicate that each network branch isolates distinct morphological features that cannot be synthesized by any single architecture in isolation:

\begin{enumerate}[leftmargin=*]
\item \textit{ConvNeXt-Base (Subchondral Trabecular Texture):} Depthwise separable convolutions operating over $7 \times 7$ receptive fields prove adept at resolving high-frequency subchondral bone eburnation and sharp marginal osteophytic spikes along the articular rim. Zeroing this branch induces a 7.32\% reduction in diagnostic accuracy, demonstrating that local trabecular density patterns provide indispensable diagnostic evidence.
\item \textit{Swin Transformer (Inter-Compartmental Relational Context):} Shifted-window self-attention captures long-range spatial dependencies across the articular gap, evaluating medial-to-lateral compartment height asymmetry and overall coronal joint alignment. Without this contextual perspective, classification accuracy drops by 4.99\%.
\item \textit{DINO ViT-Base (Invariant Articular Geometry):} Preserving frozen self-supervised weights contributes an objective structural prior that resists overfitting on modest medical imaging cohorts. Ablating DINO produces the sharpest performance deterioration ($-7.64\%$ accuracy, $-0.113$ QWK), indicating that contrastively learned anatomical boundaries serve as an essential geometric anchor for the downstream classifier.
\end{enumerate}

\subsection{Critical Methodological Constraints}
\label{subsec:critical_limitations}

A transparent assessment of methodological constraints is essential for guiding future research and clinical translation:

\subsubsection{Cohort Source and Generalizability}
Although our dataset comprises 9,786 radiographs across 4,130 unique patients, all images originate from the Osteoarthritis Initiative (OAI) cohort. While the OAI is widely recognized as a benchmark reference, its participants were recruited across four United States clinical centers under standardized radiographic acquisition protocols. Consequently, the model's robustness against domain shifts introduced by diverse radiographic detector hardware, varied exposure parameters, and international patient demographics remains to be prospectively validated.

\subsubsection{Resolution and Full-Frame Input Trade-offs}
Standardizing radiograph resolution to $224 \times 224$ pixels facilitates compatibility with pretrained vision backbones but introduces downsampling that may attenuate subtle micro-structural features, such as minimal tibial spine sharpening or faint subchondral cystic radiolucencies. While our multi-method explainability maps confirm that attention concentrates appropriately on the tibiofemoral joint space, incorporating localized high-resolution anatomical bounding boxes could further enhance discriminative precision along borderline KL transitions.

\subsubsection{Three-Tier Consolidation versus Five-Grade Granularity}
Consolidating the five KL grades into Healthy (KL~0--1), Moderate (KL~2--3), and Severe (KL~4) reflects established clinical management thresholds (observation vs.\ conservative therapy vs.\ surgical referral). However, this consolidation precludes granular discrimination between KL~2 and KL~3, which may be relevant for clinical trials monitoring slow disease progression over multi-year intervals.

\subsection{Strategic Future Horizons}
\label{subsec:future_scope}

Building upon the empirical findings and methodological constraints of this study, three high-priority research trajectories are identified:

\subsubsection{Domain-Specific Self-Supervised Foundation Pretraining}
Rather than relying on ImageNet pretraining, future work should explore pretraining vision transformers directly on large, uncurated musculoskeletal imaging archives (e.g., MIMIC-CXR, CheXpert, and large-scale extremity PACS databases) via masked autoencoding (MAE) or DINO distillation. Pretraining on radiographic physics and trabecular morphology would yield feature representations tailored to musculoskeletal disorders.

\subsubsection{Multimodal Longitudinal Progression Forecasting}
Integrating radiographic image representations with non-imaging clinical biomarkers---including age, body mass index (BMI), serum/urinary cartilage oligomeric matrix protein (COMP) biomarkers, and patient-reported Western Ontario and McMaster Universities Osteoarthritis Index (WOMAC) scores---represents a promising direction for disease forecasting. Combining temporal graph neural networks with serial radiographic examinations could enable accurate 24-month joint collapse risk prediction.

\subsubsection{Federated Multi-Center Prospective Validation}
To establish multi-center generalizability without compromising patient privacy, deploying the proposed framework within a federated learning consortium across international hospital networks represents the most viable path toward clinical adoption and regulatory clearance.


% ==============================================================================
% 6. CONCLUSION
% ==============================================================================
\section{Conclusion}
\label{sec:conclusion}

This investigation introduced an automated, leakage-free hybrid deep learning framework for knee osteoarthritis severity grading, integrating ConvNeXt-Base, Swin Transformer, and self-supervised DINO ViT through late feature fusion. Evaluated across 9,786 patient-stratified knee radiographs (4,130 unique subjects) with verified zero patient-level leakage, the framework achieved 90.92\% accuracy, a Quadratic Weighted Kappa of 0.853, and a macro ROC-AUC of 0.980 on a 4,008-image evaluation cohort under dual-pass test-time augmentation. Importantly, the Severe class attained 100\% sensitivity with zero extreme grade reversals. Explainability analysis across seven attribution methods confirmed consistent anatomical grounding on the tibiofemoral joint space. Future investigations will prioritize domain-specific pretraining on musculoskeletal radiographs, multimodal clinical feature integration, and multi-center prospective validation.


% ==============================================================================
% DATA AVAILABILITY STATEMENT
% ==============================================================================
\section*{Data Availability Statement}
The knee radiograph dataset analyzed in this study is publicly accessible via Mendeley Data (DOI: 10.17632/56rmx5bjcr.1, Version~1) and sourced from the Osteoarthritis Initiative (OAI, \url{https://oai.epi-ucsf.org/datarelease}). All code for patient-level leakage verification, model training, evaluation, and explainability is available at \url{https://github.com/SettyBhavithav/Knee-OA-Diagnostic-Pipeline-Hybrid-ViT-Ensemble}.


% ==============================================================================
% REFERENCES
% ==============================================================================
\FloatBarrier
\begin{thebibliography}{55}
\bibitem{hunter2019oa}
D. J. Hunter and S. Bierma-Zeinstra, ``Osteoarthritis,'' \emph{The Lancet}, vol. 393, no. 10182, pp. 1745--1759, 2019.

\bibitem{loeser2012oa}
R. F. Loeser, S. R. Goldring, C. R. Scanzello, and M. B. Goldring, ``Osteoarthritis: A disease of the joint as an organ,'' \emph{Arthritis \& Rheumatism}, vol. 64, no. 6, pp. 1697--1707, 2012.

\bibitem{cross2014global}
M. Cross, E. Smith, D. Hoy, S. Nolte, I. Ackerman, M. Fransen, L. Bridgett, S. Williams, F. Guillemin, C. L. Hill \emph{et al.}, ``The global burden of hip and knee osteoarthritis: Estimates from the global burden of disease 2010 study,'' \emph{Annals of the Rheumatic Diseases}, vol. 73, no. 7, pp. 1323--1330, 2014.

\bibitem{woolf2003burden}
A. D. Woolf and B. Pfleger, ``Burden of major musculoskeletal conditions,'' \emph{Bulletin of the World Health Organization}, vol. 81, no. 9, pp. 646--656, 2003.

\bibitem{braun2012diagnosis}
H. J. Braun and G. E. Gold, ``Diagnosis of osteoarthritis: Imaging,'' \emph{Bone}, vol. 51, no. 2, pp. 278--288, 2012.

\bibitem{kellgren1957radiological}
J. H. Kellgren and J. S. Lawrence, ``Radiological assessment of osteo-arthrosis,'' \emph{Annals of the Rheumatic Diseases}, vol. 16, no. 4, pp. 494--502, 1957.

\bibitem{wright2014interobserver}
R. W. Wright and MARS Group, ``Osteoarthritis classification scales: Interobserver reliability and arthroscopic correlation,'' \emph{The Journal of Bone and Joint Surgery}, vol. 96, no. 14, pp. 1145--1151, 2014.

\bibitem{sheehy2015validity}
L. Sheehy, E. Culham, L. McLean, J. Niu, J. Lynch, N. A. Segal, J. A. Singh, M. Nevitt, and D. T. Felson, ``Validity and sensitivity to change of three scales for the radiographic assessment of knee osteoarthritis using images from the multicenter osteoarthritis study (MOST),'' \emph{Osteoarthritis and Cartilage}, vol. 23, no. 9, pp. 1491--1498, 2015.

\bibitem{culvenor2019oa}
A. G. Culvenor, B. E. Oiestad, H. F. Hart, J. J. Stefanik, A. Guermazi, and K. M. Crossley, ``Prevalence of knee osteoarthritis features on magnetic resonance imaging in asymptomatic uninjured adults: A systematic review and meta-analysis,'' \emph{British Journal of Sports Medicine}, vol. 53, no. 20, pp. 1268--1278, 2019.

\bibitem{tiulpin2018automatic}
A. Tiulpin, J. Thevenot, E. Rahtu, P. Lehenkari, and S. Saarakkala, ``Automatic knee osteoarthritis diagnosis from plain radiographs: A deep learning-based approach,'' \emph{Scientific Reports}, vol. 8, no. 1, p. 1727, 2018.

\bibitem{thomas2020automated}
K. A. Thomas, L. Kidzinski, E. Halilaj, S. L. Fleming, G. R. Venkataraman, E. H. G. Oei, G. E. Gold, and S. L. Delp, ``Automated classification of radiographic knee osteoarthritis severity using deep neural networks,'' \emph{Radiology: Artificial Intelligence}, vol. 2, no. 2, p. e190065, 2020.

\bibitem{chen2019fully}
P. Chen, L. Gao, X. Shi, K. Allen, and L. Yang, ``Fully automatic knee osteoarthritis severity grading using deep neural networks with a novel ordinal loss,'' \emph{Computerized Medical Imaging and Graphics}, vol. 75, pp. 84--92, 2019.

\bibitem{pierson2021algorithmic}
E. Pierson, D. M. Cutler, J. Leskovec, S. Mullainathan, and Z. Obermeyer, ``An algorithmic approach to reducing unexplained pain disparities in underserved populations,'' \emph{Nature Medicine}, vol. 27, no. 1, pp. 136--140, 2021.

\bibitem{antony2016quantifying}
J. Antony, K. McGuinness, N. E. O'Connor, and K. Moran, ``Quantifying radiographic knee osteoarthritis severity using deep convolutional neural networks,'' in \emph{Proceedings of the 23rd International Conference on Pattern Recognition (ICPR)}. IEEE, 2016, pp. 1195--1200.

\bibitem{schiphof2008differences}
D. Schiphof, M. Boers, and S. M. A. Bierma-Zeinstra, ``Differences in descriptions of Kellgren and Lawrence grades of knee osteoarthritis,'' \emph{Annals of the Rheumatic Diseases}, vol. 67, no. 7, pp. 1034--1036, 2008.

\bibitem{zhang2010eular}
W. Zhang, M. Doherty, G. Peat, S. M. A. Bierma-Zeinstra, N. K. Arden, B. Bresnihan, G. Herrero-Beaumont, S. Kirschner, B. F. Leeb, L. S. Lohmander \emph{et al.}, ``EULAR evidence-based recommendations for the diagnosis of knee osteoarthritis,'' \emph{Annals of the Rheumatic Diseases}, vol. 69, no. 3, pp. 483--489, 2010.

\bibitem{emrani2008joint}
P. S. Emrani, J. N. Katz, C. L. Kessler, W. M. Reichmann, E. A. Wright, T. E. McAlindon, and E. Losina, ``Joint space narrowing and Kellgren--Lawrence progression in knee osteoarthritis: An analytic literature synthesis,'' \emph{Osteoarthritis and Cartilage}, vol. 16, no. 8, pp. 873--882, 2008.

\bibitem{shamir2009early}
L. Shamir, S. M. Ling, W. W. Scott, A. Bos, N. Orlov, T. J. Macura, D. M. Eckley, L. Ferrucci, and I. G. Goldberg, ``Knee x-ray image analysis method for automated detection of osteoarthritis,'' \emph{IEEE Transactions on Biomedical Engineering}, vol. 56, no. 2, pp. 407--415, 2009.

\bibitem{oka2008fully}
H. Oka, S. Muraki, T. Akune, A. Mabuchi, T. Suzuki, H. Yoshida, S. Yamamoto, K. Nakamura, N. Yoshimura, and H. Kawaguchi, ``Fully automatic quantification of knee osteoarthritis severity on plain radiographs,'' \emph{Osteoarthritis and Cartilage}, vol. 16, no. 11, pp. 1300--1306, 2008.

\bibitem{thomson2015automated}
J. Thomson, T. O'Neill, D. Felson, and T. Cootes, ``Automated shape and texture analysis for detection of osteoarthritis from radiographs of the knee,'' \emph{Medical Image Computing and Computer-Assisted Intervention (MICCAI)}, 2015.

\bibitem{rajpurkar2017chexnet}
P. Rajpurkar, J. Irvin, K. Zhu, B. Yang, H. Mehta, T. Duan, D. Ding, A. Bagul, C. Langlotz, K. Shpanskaya \emph{et al.}, ``CheXNet: Radiologist-level pneumonia detection on chest x-rays with deep learning,'' \emph{arXiv preprint arXiv:1711.05225}, 2017.

\bibitem{gulshan2016development}
V. Gulshan, L. Peng, M. Coram, M. C. Stumpe, D. Wu, A. Narayanaswamy, S. Venugopalan, K. Widner, T. Madams, J. Cuadros \emph{et al.}, ``Development and validation of a deep learning algorithm for detection of diabetic retinopathy in retinal fundus photographs,'' \emph{JAMA}, vol. 316, no. 22, pp. 2402--2410, 2016.

\bibitem{bejnordi2017diagnostic}
B. Ehteshami Bejnordi, M. Veta, P. J. van Diest, B. van Ginneken, N. Karssemeijer, G. Litjens, J. A. W. M. van der Laak, and the CAMELYON16 Consortium, ``Diagnostic assessment of deep learning algorithms for detection of lymph node metastases in women with breast cancer,'' \emph{JAMA}, vol. 318, no. 22, pp. 2199--2210, 2017.

\bibitem{he2016deep}
K. He, X. Zhang, S. Ren, and J. Sun, ``Deep residual learning for image recognition,'' in \emph{Proceedings of the IEEE Conference on Computer Vision and Pattern Recognition (CVPR)}, 2016, pp. 770--778.

\bibitem{huang2017densely}
G. Huang, Z. Liu, L. van der Maaten, and K. Q. Weinberger, ``Densely connected convolutional networks,'' in \emph{Proceedings of the IEEE Conference on Computer Vision and Pattern Recognition (CVPR)}, 2017, pp. 4700--4708.

\bibitem{tan2019efficientnet}
M. Tan and Q. V. Le, ``EfficientNet: Rethinking model scaling for convolutional neural networks,'' in \emph{International Conference on Machine Learning (ICML)}. PMLR, 2019, pp. 6105--6114.

\bibitem{liu2021swin}
Z. Liu, Y. Lin, Y. Cao, H. Hu, Y. Wei, Z. Zhang, S. Lin, and B. Guo, ``Swin Transformer: Hierarchical vision transformer using shifted windows,'' in \emph{Proceedings of the IEEE/CVF International Conference on Computer Vision (ICCV)}, 2021, pp. 10012--10022.

\bibitem{liu2022convnext}
Z. Liu, H. Mao, C.-Y. Wu, C. Feichtenhofer, T. Darrell, and S. Xie, ``A ConvNet for the 2020s,'' in \emph{Proceedings of the IEEE/CVF Conference on Computer Vision and Pattern Recognition (CVPR)}, 2022, pp. 11976--11986.

\bibitem{caron2021dino}
M. Caron, H. Touvron, I. Misra, H. Jegou, J. Mairal, P. Bojanowski, and A. Joulin, ``Emerging properties in self-supervised vision transformers,'' in \emph{Proceedings of the IEEE/CVF International Conference on Computer Vision (ICCV)}, 2021, pp. 9650--9660.

\bibitem{raghu2019transfusion}
M. Raghu, C. Zhang, J. Kleinberg, and S. Bengio, ``Transfusion: Understanding transfer learning for medical imaging,'' \emph{Advances in Neural Information Processing Systems (NeurIPS)}, vol. 32, 2019.

\bibitem{tajbakhsh2016convolutional}
N. Tajbakhsh, J. Y. Shin, S. R. Gurudu, R. T. Hurst, C. B. Kendall, M. B. Gotway, and J. Liang, ``Convolutional neural networks for medical image analysis: Full training or fine tuning?'' \emph{IEEE Transactions on Medical Imaging}, vol. 35, no. 5, pp. 1299--1312, 2016.

\bibitem{antony2017automatic}
J. Antony, K. McGuinness, K. Moran, and N. E. O'Connor, ``Automatic detection of knee joints and quantification of knee osteoarthritis severity using convolutional neural networks,'' in \emph{Machine Learning and Data Mining in Pattern Recognition (MLDM)}. Springer, 2017, pp. 376--390.

\bibitem{swiecicki2021deep}
A. Swiecicki, D. Kondziella, O. Puonti, H. R. Siebner, and M. Blinkenberg, ``Deep learning for automatic grading of knee osteoarthritis on radiographs,'' \emph{Osteoarthritis and Cartilage}, 2021.

\bibitem{guan2022deep}
B. Guan, F. Liu, A. Hasan, H. Bilal, and M. Shafiq, ``Deep learning approach to predict radiographic knee osteoarthritis,'' \emph{Journal of Healthcare Engineering}, vol. 2022, pp. 1--12, 2022.

\bibitem{bany2023knee}
M. Bany Muhammad and A. S. Tarawneh, ``Knee osteoarthritis detection and severity classification using residual neural networks on preprocessed x-ray images,'' \emph{Diagnostics}, vol. 13, no. 8, p. 1380, 2023.

\bibitem{kondal2023transfer}
S. R. Kondal, A. Kulkarni \emph{et al.}, ``Transfer learning-based knee osteoarthritis grading using vision transformers,'' \emph{Biomedical Signal Processing and Control}, 2023.

\bibitem{bayramoglu2021adaptive}
N. Bayramoglu, A. Tiulpin, J. Hirvasniemi, M. Niemelainen, and S. Saarakkala, ``Adaptive segmentation of knee radiographs for selecting the optimal pretrained model for severity assessment of osteoarthritis,'' \emph{Computerized Medical Imaging and Graphics}, vol. 88, p. 101831, 2021.

\bibitem{chen2018dataset}
P. Chen, ``Knee osteoarthritis severity grading dataset,'' \emph{Mendeley Data}, vol. V1, 2018.

\bibitem{carr2012knee}
A. J. Carr, O. Robertsson, S. Graves, A. J. Price, N. K. Arden, A. Judge, and D. J. Beard, ``Knee replacement,'' \emph{The Lancet}, vol. 379, no. 9823, pp. 1331--1340, 2012.

\bibitem{zhang2018mixup}
H. Zhang, M. Cisse, Y. N. Dauphin, and D. Lopez-Paz, ``mixup: Beyond empirical risk minimization,'' \emph{International Conference on Learning Representations (ICLR)}, 2018.

\bibitem{thulasidasan2019mixup}
S. Thulasidasan, G. Chennupati, J. A. Bilmes, T. Bhattacharya, and S. Michalak, ``On mixup training: Improved calibration and predictive uncertainty for deep neural networks,'' \emph{Advances in Neural Information Processing Systems (NeurIPS)}, vol. 32, 2019.

\bibitem{loshchilov2019adamw}
I. Loshchilov and F. Hutter, ``Decoupled weight decay regularization,'' \emph{International Conference on Learning Representations (ICLR)}, 2019.

\bibitem{loshchilov2017sgdr}
I. Loshchilov and F. Hutter, ``SGDR: Stochastic gradient descent with warm restarts,'' \emph{International Conference on Learning Representations (ICLR)}, 2017.

\bibitem{landis1977measurement}
J. R. Landis and G. G. Koch, ``The measurement of observer agreement for categorical data,'' \emph{Biometrics}, vol. 33, no. 1, pp. 159--174, 1977.

\bibitem{reyes2020interpretability}
M. Reyes, R. Meier, S. Pereira, C. A. Silva, F.-M. Dahlweid, H. von Tengg-Kobligk, R. M. Summers, and R. Wiest, ``On the interpretability of artificial intelligence in radiology: Challenges and opportunities,'' \emph{Radiology: Artificial Intelligence}, vol. 2, no. 3, p. e190043, 2020.

\bibitem{selvaraju2017gradcam}
R. R. Selvaraju, M. Cogswell, A. Das, R. Vedantam, D. Parikh, and D. Batra, ``Grad-CAM: Visual explanations from deep networks via gradient-based localization,'' in \emph{Proceedings of the IEEE International Conference on Computer Vision (ICCV)}, 2017, pp. 618--626.

\bibitem{chattopadhay2018gradcampp}
A. Chattopadhay, A. Sarkar, P. Howlader, and V. N. Balasubramanian, ``Grad-CAM++: Generalized gradient-based visual explanations for deep convolutional networks,'' in \emph{IEEE Winter Conference on Applications of Computer Vision (WACV)}, 2018, pp. 839--847.

\bibitem{jiang2021layercam}
P.-T. Jiang, C.-B. Zhang, Q. Hou, M.-M. Cheng, and Y. Wei, ``LayerCAM: Exploring hierarchical class activation maps for localization,'' \emph{IEEE Transactions on Image Processing}, vol. 30, pp. 5875--5888, 2021.

\bibitem{wang2020scorecam}
H. Wang, Z. Wang, M. Du, F. Yang, Z. Zhang, S. Ding, P. Mardziel, and X. Hu, ``Score-CAM: Score-weighted visual explanations for convolutional neural networks,'' in \emph{Proceedings of the IEEE/CVF Conference on Computer Vision and Pattern Recognition Workshops (CVPRW)}, 2020, pp. 24--25.

\bibitem{sundararajan2017axiomatic}
M. Sundararajan, A. Taly, and Q. Yan, ``Axiomatic attribution for deep networks,'' in \emph{International Conference on Machine Learning (ICML)}. PMLR, 2017, pp. 3319--3328.

\bibitem{ribeiro2016lime}
M. T. Ribeiro, S. Singh, and C. Guestrin, ````Why Should I Trust You?'': Explaining the predictions of any classifier,'' in \emph{Proceedings of the 22nd ACM SIGKDD International Conference on Knowledge Discovery and Data Mining}. ACM, 2016, pp. 1135--1144.

\bibitem{johnson2019survey}
J. M. Johnson and T. M. Khoshgoftaar, ``Survey on deep learning with class imbalance,'' \emph{Journal of Big Data}, vol. 6, no. 1, p. 27, 2019.

\bibitem{howard2018ulmfit}
J. Howard and S. Ruder, ``Universal language model fine-tuning for text classification,'' in \emph{Proceedings of the 56th Annual Meeting of the Association for Computational Linguistics (ACL)}, 2018, pp. 328--339.

\bibitem{shanmugam2021tta}
D. Shanmugam, D. Blalock, G. Balakrishnan, and J. Guttag, ``Better aggregation in test-time augmentation,'' \emph{Proceedings of the IEEE/CVF International Conference on Computer Vision (ICCV)}, pp. 1214--1223, 2021.

\bibitem{nevitt2006oai}
M. C. Nevitt, D. T. Felson, and G. Lester, ``The osteoarthritis initiative: Protocol for the cohort study,'' 2006, available at: https://oai.epi-ucsf.org.
\end{thebibliography}

\end{document}
